% 87-02 ②(87쪽) — 보기 그래프: $x=-1$ 에서 기울기 0, $x=3$ 에서 극대인 위로 볼록한 사차함수 개형. 그림 자체가 보기라 TikZ 로 다시 그림.
% 라벨은 원문 그대로 -1, O, 3, x, y, y=f(x) 와 점선 두 줄만.
\begin{tikzpicture}[baseline=(current bounding box.center), x=0.36cm, y=0.36cm, line width=0.5pt]
  \path (-2.9,-2.4) rectangle (5.15,3.5);
  \draw[->, >=stealth, line width=0.7pt] (-2.6,0) -- (4.85,0) node[below=1pt, font=\footnotesize] {$x$};
  \draw[->, >=stealth, line width=0.7pt] (0,-2.1) -- (0,3.25) node[left=1pt, font=\footnotesize] {$y$};
  \draw[densely dotted, line width=0.4pt] (-1,0) -- (-1,0.3);
  \draw[densely dotted, line width=0.4pt] (3,0) -- (3,2.6);
  \draw[line width=0.6pt, smooth] plot coordinates {(-2.3,-1.7) (-2.15,-1.1) (-2,-0.6) (-1.8,-0.15) (-1.55,0.12) (-1.3,0.25) (-1,0.3) (-0.7,0.42) (-0.4,0.62) (0,0.95) (0.5,1.4) (1,1.8) (1.5,2.15) (2,2.4) (2.5,2.55) (2.8,2.6) (3,2.6) (3.2,2.5) (3.45,2.25) (3.7,1.7) (3.9,1.0) (4.05,0.2) (4.15,-0.6) (4.25,-1.5)};
  \node[below=1pt, font=\footnotesize] at (-1.3,0) {$-1$};
  \node[below left=-2pt and -3pt, font=\footnotesize] at (0,0) {$\pt{O}$};
  \node[below=1pt, font=\footnotesize] at (3,0) {$3$};
  \node[anchor=west, font=\footnotesize] at (0.25,3.0) {$y=f(x)$};
\end{tikzpicture}
